\documentclass[11pt,a4paper,fleqn]{article}
\usepackage[T1]{fontenc}
\usepackage[utf8]{inputenc}
\usepackage{lmodern}
\usepackage[czech]{babel}
\usepackage[margin=1.7cm]{geometry}
\usepackage{amsmath,amssymb,mathtools}
\usepackage{xcolor,cancel,enumitem}
\usepackage{tikz,pgfplots}
\pgfplotsset{compat=1.18}
\usetikzlibrary{arrows.meta,calc,patterns,shapes.geometric,babel}
\setlength{\mathindent}{1.5em}
% červeně zakroužkovaný a přeškrtnutý výraz
\newcommand{\kr}[1]{\mathpalette\krA{#1}}
\newcommand{\krA}[2]{\tikz[baseline=(k.base)]\node[draw=red!75!black, ellipse, inner sep=0pt] (k) {$#1#2$};}
\renewcommand{\CancelColor}{\color{red!75!black}}
\newcommand{\cerv}[1]{{\color{red!75!black}#1}}
% značka pro šipky (2 průchody pdflatex)
\newcommand{\tm}[2]{\tikz[remember picture, baseline=(#1.base)]\node[inner sep=0pt, outer sep=1pt] (#1) {$\displaystyle #2$};}
% malý graf ln x
\newcommand{\lngraf}{\begin{tikzpicture}[scale=0.45]
\draw[-{Stealth}] (-0.6,0) -- (2.8,0);
\draw[-{Stealth}] (0,-1.4) -- (0,1.4);
\draw[thick, domain=0.25:2.7, samples=40] plot (\x, {ln(\x)}) node[right, font=\scriptsize] {$\ln$};
\end{tikzpicture}}
\pagestyle{empty}
\setlength{\parindent}{0pt}
\setlength{\parskip}{4pt}
\newcommand{\vysl}[1]{\boxed{\color{red!75!black}#1}}
\newcommand{\pozn}[1]{{\color{red!75!black}#1}}
\newcommand{\dx}{\,\mathrm{d}x}
\newcommand{\dt}{\,\mathrm{d}t}
\newcommand{\subst}[1]{\left|\begin{array}{l}#1\end{array}\right|}
\DeclareMathOperator{\arctg}{arctg}
\DeclareMathOperator{\arccotg}{arccotg}
\DeclareMathOperator{\tg}{tg}
\DeclareMathOperator{\cotg}{cotg}
\usepackage{tcolorbox}
\newtcolorbox{pomocny}{colframe=red!75!black, colback=white, coltitle=white, fonttitle=\bfseries, title={Pomocný výpočet:}, sharp corners}
\newcommand{\dy}{\,\mathrm{d}y}

\newcommand{\priblizeni}[3]{%
\begin{tikzpicture}[baseline=-3pt]
\draw (-1.5,0) -- (1.5,0);
\draw (0,0.1) -- (0,-0.1) node[below] {$#1$};
\draw[red!75!black, -{Stealth}] (1.2,-0.3) -- (0.3,-0.3) node[midway, above] {\small $#2$};
\node[red!75!black, below] at (0.9,-0.35) {\small $#3$};
\end{tikzpicture}}

\begin{document}

\section*{Z grafu funkce určete limity}

\textbf{a)}\quad
$\displaystyle\lim_{x\to\infty}e^x=\vysl{\infty}$
\qquad
$\displaystyle\lim_{x\to-\infty}e^x=\vysl{0}$
\qquad
$\displaystyle\lim_{x\to0}e^x=\vysl{1}$
\qquad
\pozn{$\Big($}$\displaystyle\lim_{x\to1}e^x=\vysl{e}$\pozn{$\Big)$ lze počítat i dosazením}

\begin{center}
\begin{tikzpicture}
\begin{axis}[axis lines=middle, width=6cm, height=4.5cm, clip=false,
  xmin=-2.5, xmax=2, ymin=-0.8, ymax=4.5, xtick={1}, ytick={1,2.71828}, yticklabels={$1$,$e$}]
\addplot[thick, domain=-2.5:1.5, samples=60] {exp(x)} node[right] {$e^x$};
\addplot[only marks, mark=*] coordinates {(0,1)};
\draw[dotted] (axis cs:0,2.71828) -- (axis cs:1,2.71828) -- (axis cs:1,0);
\end{axis}
\end{tikzpicture}
\end{center}

\textbf{b)}\quad
$\displaystyle\lim_{x\to\infty}\arccotg x=\vysl{0}$
\qquad
$\displaystyle\lim_{x\to-\infty}\arccotg x=\vysl{\pi}$

\begin{minipage}{0.5\textwidth}
\pozn{$\Big($}$\displaystyle\lim_{x\to1}\arccotg x=\arccotg1=\pozn{45^\circ=\frac\pi4}$\pozn{$\Big)$}\\[4pt]
\begin{tikzpicture}[scale=0.85]
\node[red!75!black, align=left, anchor=east] at (-1.6,-0.6) {určíme\\ z jednotkové kružnice};
\draw (-1.4,0) -- (1.5,0);
\draw (0,-1.3) -- (0,1.5);
\draw (0,0) circle (1);
\draw (-1.4,1) -- (1.7,1) node[right] {\small $\arccotg$};
\draw (1,1.5) -- (1,-1.3);
\fill (0,1) circle (1.5pt);
\fill (1,1) circle (1.8pt) node[above right=-1pt] {\small $1$};
\draw[red!75!black, dashed] (0,0) -- (1,1);
\fill[red!75!black] (45:1) circle (1.8pt);
\draw[red!75!black] (0.35,0) arc (0:45:0.35);
\node[red!75!black, font=\scriptsize] at (18:0.62) {$45^\circ$};
\end{tikzpicture}
\end{minipage}%
\begin{minipage}{0.5\textwidth}
\begin{tikzpicture}
\begin{axis}[axis lines=middle, width=8cm, height=4.5cm, clip=false,
  xmin=-6, xmax=6, ymin=-0.6, ymax=3.8, xtick=\empty, ytick={1.5708,3.1416}, yticklabels={$\frac{\pi}{2}$,$\pi$}]
\addplot[dashed, domain=-6:6] {pi};
\addplot[thick, domain=-6:6, samples=100] {pi/2 - rad(atan(x))} node[right] {$\arccotg$};
\addplot[only marks, mark=*] coordinates {(0,1.5708)};
\end{axis}
\end{tikzpicture}
\end{minipage}

\section*{Vypočti:}

\textbf{a)}\quad $\displaystyle
\lim_{x\to\infty}e^{\tfrac{2x}{x^2-1}} = e^{\lim\limits_{x\to\infty}\tfrac{\kr{2x}}{\kr{x^2}-1}} = e^{\lim\limits_{x\to\infty}\tfrac{2\cancel{x}}{x^{\cancel{2}}}}
= e^{\lim\limits_{x\to\infty}\tfrac2x} = e^0 = \vysl{1}
$


\textbf{b)}\quad $\displaystyle
\lim_{x\to\infty}\frac{\kr{3x^4}-2}{\kr{2x^4}+x-3} = \lim_{x\to\infty}\frac{3\cancel{x^4}}{2\cancel{x^4}} = \vysl{\frac32}
$


\textbf{c)}
\[
\lim_{x\to-3^+}\frac{2-x}{\tm{den}{\cerv{2}x^2+5x-3}} = \left[\frac50\right] = \lim_{x\to-3^+}\frac{2-x}{\cerv{2}(x+3)\left(x-\frac12\right)}
= \lim_{x\to-3^+}\frac{1}{x+3}\cdot\lim_{x\to-3^+}\frac{2-x}{2\left(x-\frac12\right)} =
\]
\begin{minipage}{0.55\textwidth}
\begin{align*}
D &= \tm{Dl}{25-4\cdot2\cdot(-3)} = 49\\
x_{1,2} &= \frac{-5\pm\sqrt{49}}{4} = \frac{-5\pm7}{4} = \begin{cases}-3\\ \frac24=\frac12\end{cases}
\end{align*}
\end{minipage}%
\begin{minipage}{0.4\textwidth}
\priblizeni{-3}{+}{(x+3)>0}
\end{minipage}
\begin{tikzpicture}[remember picture, overlay]
\draw[-{Stealth}] ([yshift=-2pt]den.south) -- ([yshift=3pt]Dl.north);
\end{tikzpicture}
\[
= \frac{1}{0^+}\cdot\frac{5}{2\left(-3-\frac12\right)} = \infty\cdot\frac{5}{2\cdot\left(-\frac72\right)} = \infty\cdot\left(-\frac57\right) = \vysl{-\infty}
\]

\textbf{d)}
\begin{align*}
\lim_{x\to4^+}\ln\frac{3x+1}{x^2-4x} &= \ln\lim_{x\to4^+}\frac{3x+1}{x^2-4x} = \ln\lim_{x\to4^+}\frac{3x+1}{x(x-4)}
= \ln\left(\lim_{x\to4^+}\frac{1}{x-4}\cdot\lim_{x\to4^+}\frac{3x+1}{x}\right) =\\
&= \ln\left(\frac{1}{0^+}\cdot\frac{3\cdot4+1}{4}\right) = \ln\left(\infty\cdot\frac{13}{4}\right) = \tm{lninf}{\ln\infty} = \vysl{\infty}
\qquad \priblizeni{4}{}{(x-4)>0}
\end{align*}
\begin{tikzpicture}[remember picture, overlay]
\node[anchor=north, inner sep=0pt] at ([yshift=-4pt]lninf.south) {\lngraf};
\end{tikzpicture}
\vspace{0.9cm}


\textbf{e)}\quad $\displaystyle
\lim_{x\to-\infty}\left(\kr{x^3}-x+2\right) = (-\infty)^3 = \vysl{-\infty}
$

\end{document}
